\scanpage{80}
\deleted{III) Notons maintenant que dans les trois cas I II III envisagés, faisant la construction sur $\mathbb R$ et $\mathbb C$, on trouve des cartes cellulaires régulières sur la sphère -- qui est \uline{simplement connexe}. On en déduit \inserted{par la théorie des cartes} que le groupe $G_0^+$ \inserted{$(=\mathfrak A_4,\mathfrak S_4,\mathfrak A_5)$} est engendré par les gén. $\rho_s,\rho_f$ avec les seules relations $\rho_f^{\nu_0}=\rho_s^{\nu_1}=(\rho_f\rho_s)^2=1$. La théorie}

% slanted
\margin{$\bigcirc$ on pourrait sans doute l'obtenir \unsure{directement}}

\deleted{\unsure{au cas} \illegible{5} II, III,} du groupe fondamental de la sphère moins trois pts, \deleted{\unsure{IV A la} \illegible{2} I, \unsure{ou}}

\deleted{III} montre alors que dans le cas général, le ss-groupe fini $G_0$ \deleted{\unsure{engendré}} doit être un \uline{quotient} d'un des gpes $\mathfrak A_4,\mathfrak S_4,\mathfrak A_5$ -- ceci implique déjà que les ordres de ses éléments sont $\leq5$, donc

\[
 \boxed{3\leq\nu_0,\nu_1\leq5}
\]

D'ailleurs, exprimant (sur $k_0=\mathbb R$) que la forme quadratique de matrice (matrice par r. à $a_0,a_1,a_2$)

\[
 \begin{pmatrix}
 2 & -\deleted{\unsure{2}}\beta_0 & 0\\
 -\deleted{\unsure{2}}\beta_0 & 2\beta_0 & -\deleted{\unsure{2}}\beta_0\beta_1\\
 0 & -\deleted{\unsure{2}}\beta_0\beta_1 & 2\beta_0\beta_1
 \end{pmatrix}
\]

est définie positive, on trouve les conditions

\[
 \deleted{\unsure{\beta_0>0,\beta_1>0|\beta_0,\beta_1>0}}\qquad
 4-(\beta_0+\beta_1)>0\qquad\beta_0(4-\beta_0)>0
\]

\[
 \text{i.e}\qquad\left\{\begin{aligned}
 &\alpha_0+\alpha_1<0\\
 &\deleted{\unsure{\alpha_0^2<4}}\ \text{\unsure{i.e.}}\quad -2<\alpha_0,\alpha_1\unsure{<|\leq}+2
 \end{aligned}\right.
\]

(Conditions néc. et suff pour que la forme \unsure{envisagée|env. géom} soit définie, pour des \deleted{\unsure{couples}} $\alpha_0,\alpha_1$ réels).
% archive sheet number: 79
